\section{Hipoteza}
\label{sec:hipoteza}

Hipoteza ovog rada jest da se postojeći percepcijski cjevovod za detekciju i praćenje grane, razvijen za horizontalnu konfiguraciju kamere i prilazak s vanjskom lokalizacijom, može ponovno iskoristiti gotovo bez izmjena za scenarij u kojem je kamera okrenuta prema gore i u kojem letjelica ne raspolaže sustavom apsolutne lokalizacije, uz uvjet da se zamijeni isključivo upravljački sloj koji povezuje percepciju s naredbama letjelici.

\section{Ciljevi rada}
\label{sec:ciljevi}

Glavni ciljevi ovog rada su:

\begin{enumerate}
  \item Prilagoditi postojeći cjevovod za detekciju grane i praćenje točke prijanjanja (razvijen u sklopu diplomskog rada) novoj orijentaciji kamere (prema gore umjesto prema naprijed) i provjeriti da percepcijski dio sustava ostaje ispravan bez izmjena u geometriji obrade slike.
  \item Zamijeniti upravljački ROS čvor koji je ovisio o sustavu OptiTrack i o senzoru udaljenosti (ToF) za okidanje manevra prijanjanja paketom \texttt{ibvs\_perching}, koji upravlja letjelicom izravno preko MAVROS-a bez vanjske lokalizacije.
  \item Odrediti i implementirati preslikavanje osi između izlaza percepcijskog cjevovoda (pogreška u slikovnoj ravnini) i naredbi tijela letjelice (kutne brzine, brzina penjanja), prilagođeno geometriji u kojoj letjelica polazi ispod grane i prijanjanje izvodi pretežno vertikalnim približavanjem.
  \item Provesti eksperimentalnu provjeru prilagođenog sustava i usporediti ponašanje s onim opisanim u izvornom diplomskom radu.
\end{enumerate}
